\section*{अध्याय \ref{ch:b3:rovibrational-spectroscopy} --- घूर्णी और कंपनिक स्पेक्ट्रमिकी}

\begin{solution}{exo:b3:rovibrational-spectroscopy:1}
$I = h/(8\pi^2cB) = 6.62607 \times 10^{-34}/(8\pi^2 \times 2.99792 \times 10^{10} \times
1.9225) = \qty{1.4561e-46}{kg.m^2}$। $\mu = \qty{6.85621}{u} = \qty{1.13850e-26}{kg}$, अतः
$r_0 = \sqrt{I/\mu} = \qty{113.09}{pm}$।
\end{solution}

\begin{solution}{exo:b3:rovibrational-spectroscopy:2}
$2B_0(J+1)$: 20.88, 41.76, 62.64, \qty{83.52}{cm^{-1}}। $J_{\max} = \sqrt{kT/2hcB} -
\frac12 = \sqrt{208.5/20.88} - 0.5 = 2.7$: स्तर $J = 3$ सबसे अधिक
समष्टित है, और सबसे प्रबल रेखा $4 \leftarrow 3$ है, \qty{83.5}{cm^{-1}} पर।
\end{solution}

\begin{solution}{exo:b3:rovibrational-spectroscopy:3}
सूक्ष्मतरंग अवशोषण के लिए स्थायी द्विध्रुव चाहिए: केवल \ce{HCl} और \ce{H2O}।
घूर्णी रामन के लिए विषमदैशिक ध्रुवणीयता चाहिए: \ce{H2}, \ce{HCl}, \ce{N2},
\ce{CO2}, \ce{H2O}; \omterm{def:b3:rovibrational-spectroscopy:rotor-types}{गोलीय लट्टू} \ce{CH4} और \ce{SF6} नहीं।
\end{solution}

\begin{solution}{exo:b3:rovibrational-spectroscopy:4}
$N_1/N_0 = \eu^{-hc\tilde\nu/kT}$: \qty{300}{K} पर $\eu^{-13.84} = \num{9.8e-7}$;
\qty{1000}{K} पर $\eu^{-4.15} = 0.016$। \omterm{def:b3:rovibrational-spectroscopy:anharmonic}{तप्त बैंड} ध्यान देने योग्य (\omterm{def:b3:rovibrational-spectroscopy:anharmonic}{मूल संक्रमण} का
एक प्रतिशत) केवल \qty{1000}{K} के निकट होता है।
\end{solution}

\begin{solution}{exo:b3:rovibrational-spectroscopy:5}
$\tilde\omega_e - 2\tilde\omega_ex_e = 2885.3$ और $2\tilde\omega_e - 6\tilde\omega_ex_e =
5665.0$: दूसरे में से पहले का दुगुना घटाने पर $-2\tilde\omega_ex_e = -105.6$,
अतः $\tilde\omega_ex_e = \qty{52.8}{cm^{-1}}$ और $\tilde\omega_e = \qty{2990.9}{cm^{-1}}$।
\end{solution}

\begin{solution}{exo:b3:rovibrational-spectroscopy:6}\emergencystretch=3em
$D_e = 2990.95^2/(4 \times 52.82) = \qty{42341}{cm^{-1}}$; $G(0) = 1495.47 - 13.20 =
\qty{1482.3}{cm^{-1}}$; $D_0 = \qty{40859}{cm^{-1}} = \qty{488.8}{kJ/mol}$, वास्तविक
\qty{35760}{cm^{-1}} $= \qty{427.8}{kJ/mol}$ के विरुद्ध: मोर्स मॉडल 14\,\% अधिक आँकता है।
\end{solution}

\begin{solution}{exo:b3:rovibrational-spectroscopy:7}
$R(0) - P(2) = 6B_0 = 62.64$: $B_0 = \qty{10.440}{cm^{-1}}$। $R(1) - P(1) = 6B_1 = 60.80$:
$B_1 = \qty{10.133}{cm^{-1}}$। $\alpha_e = \qty{0.307}{cm^{-1}}$। $R(0) = \tilde\nu_0 + 2B_1$ से
$\tilde\nu_0 = \qty{2885.31}{cm^{-1}}$ मिलता है (और $P(1) = \tilde\nu_0 - 2B_0$ मेल खाता है)।
\end{solution}

\begin{solution}{exo:b3:rovibrational-spectroscopy:8}
$\mu_H = \qty{0.97959}{u}$, $\mu_D = \qty{1.90441}{u}$, $\mu_H/\mu_D = 0.51438$। $B_0(\ce{DCl})
= 10.440 \times 0.51438 = \qty{5.370}{cm^{-1}}$। $\tilde\omega_e = 2990.95 \times \sqrt{0.51438}
= 2145.1$, $\tilde\omega_ex_e = 52.82 \times 0.51438 = 27.17$: $\tilde\nu_0 = 2145.1 - 54.3 =
\qty{2090.8}{cm^{-1}}$।
\end{solution}

\begin{solution}{exo:b3:rovibrational-spectroscopy:9}
विस्थापन $4B_0(J + \frac32)$: 11.94, 19.90, \qty{27.85}{cm^{-1}} ($J = 0$, 1, 2 से)। दोनों
\ce{^{14}N} नाभिक चक्रण 1 वाले सर्वसम बोसॉन हैं: सम-$J$ स्तरों की छह
नाभिकीय-चक्रण अवस्थाएँ हैं, विषम-$J$ की तीन, अतः सम $J$ से आने वाली रेखाएँ अपने
पड़ोसियों से दुगुनी प्रबल हैं।
\end{solution}

\begin{solution}{exo:b3:rovibrational-spectroscopy:10}
\omterm{def:b3:quantum-model-systems:rotor}{दृढ़ घूर्णक} रेखाओं को $\nu_1$, $2\nu_1$, $3\nu_1$ पर रखता: $230542.40$ और
\qty{345813.61}{MHz}, जो मापी गई रेखाओं से 4.4 और \qty{17.6}{MHz} ऊपर हैं।
$\nu_{J+1\leftarrow J} = 2B(J+1) - 4D(J+1)^3$ के साथ: $\nu_1 = 2B - 4D$, $\nu_3 = 6B - 108D$, अतः
$3\nu_1 - \nu_3 = 96D$: $D = \qty{0.1835}{MHz}$ और $B = (\nu_1 + 4D)/2 =
\qty{57635.97}{MHz}$। जाँच: $\nu_2 = 4B - 32D = \qty{230538.0}{MHz}$।
\end{solution}

\begin{solution}{exo:b3:rovibrational-spectroscopy:11}
$\Delta G(v + \frac12) = \tilde\omega_e - 2\tilde\omega_ex_e(v+1)$ लगभग $v =
\tilde\omega_e/2\tilde\omega_ex_e - 1$ पर लुप्त होता है; $n$ पदों की समांतर श्रेढ़ी का योग, जो $\tilde\omega_e - 2\tilde\omega_ex_e$ से लगभग 0 तक
घटती है, लगभग $n\tilde\omega_e/2 \approx
\tilde\omega_e^2/4\tilde\omega_ex_e - \tilde\omega_e/2 \approx D_e - G(0)$ है। \ce{HCl} के लिए 28 अंतरालों
($v = 0$ से 27) का योग \qty{40857}{cm^{-1}} है, $D_e - G(0) = \qty{40859}{cm^{-1}}$ के विरुद्ध।
\end{solution}

\begin{solution}{exo:b3:rovibrational-spectroscopy:12}
केवल सम $J$ होने पर संक्रमण $J \to J + 2$, $J = 0, 2, 4, \dots$ से आरंभ होते हैं: विस्थापन
$6B, 14B, 22B, \dots$, अंतराल $8B$। \ce{CO2} में सममिति केंद्र है और कोई
स्थायी द्विध्रुव नहीं: कोई सूक्ष्मतरंग स्पेक्ट्रम नहीं।
\end{solution}

\begin{solution}{pb:b3:rovibrational-spectroscopy:1}
\textbf{1.} $B_0 = \nu/2 = \qty{57.6356}{GHz}$; $B_0/c = \qty{1.92252}{cm^{-1}}$।
\textbf{2.} $\mu = 12 \times 15.99491/27.99491 = \qty{6.85621}{u} = \qty{1.13850e-26}{kg}$।
\textbf{3.} $I = h/8\pi^2B_0 = \qty{1.4560e-46}{kg.m^2}$।
\textbf{4.} $r_0 = \sqrt{I/\mu} = \qty{113.09}{pm}$।
\textbf{5.} $r_0$, $r_e = \qty{112.83}{pm}$ से \qty{0.26}{pm} अधिक है: $B_0$, अनावर्ती कूप में
शून्य-बिंदु कंपन पर $1/r^2$ का औसत लेता है, और वह कूप लंबी दूरियों पर अधिक समय
बिताता है।
\textbf{6.} $2\nu_{1\leftarrow0} = \qty{230.5424}{GHz}$; मापी गई रेखा \qty{4.4}{MHz} नीचे है:
अपकेंद्री विरूपण।
\textbf{7.} $0$, 3.845 और \qty{11.535}{cm^{-1}}, अर्थात् $E/k = 0$, 5.53 और \qty{16.60}{K}।
\textbf{8.} $N_2/N_1 = \frac53\exp[-(16.60 - 5.53)\,\mathrm K/T]$।
\textbf{9.} $\ln(\frac53/0.85) = 0.673 = 11.06/T$: $T = \qty{16}{K}$।
\textbf{10.} $J_{\max} = \sqrt{16.4/(2 \times 1.43878 \times 1.9225)} - 0.5 = 1.2$: $J = 1$।
\textbf{11.} एक कंपनिक क्वांटम $hc\tilde\nu/k \approx \qty{3080}{K}$ के संगत है:
\qty{16}{K} पर कोई अणु कंपनिक रूप से उत्तेजित नहीं होता, अतः मेघ कोई कंपनिक
प्रकाश उत्सर्जित नहीं करता।
\textbf{12.} उसे केवल \qty{5.5}{K} के उत्तेजन की आवश्यकता है, वह सबसे ठंडी
गैस में भी उत्तेजित हो जाती है, और CO प्रचुर है तथा उसका द्विध्रुव आघूर्ण है।
\textbf{13.} $\mu = 13.00335 \times 15.99491/28.99826 = \qty{7.17241}{u}$।
\textbf{14.} $B \propto 1/\mu$: $115.2712 \times 6.85621/7.17241 = \qty{110.1893}{GHz}$।
\textbf{15.} भविष्यवाणी मापी गई रेखा से \qty{12}{MHz} (0.011\,\%) नीचे है: $B_0 = B_e - \frac12\alpha_e$, और
कंपन--घूर्णन पद $\alpha_e$, $\mu$ की किसी भिन्न घात के साथ बदलता है;
साम्य संरचना वही है, शून्य-बिंदु औसत नहीं।
\textbf{16.} $98.9/1.1 = 90$।
\textbf{17.} \ce{^{12}C^{16}O} रेखा संतृप्त (प्रकाशतः मोटी) है: वह अब
गैस की मात्रा के साथ नहीं बढ़ती।
\textbf{18.} दुर्लभ \ce{^{13}C^{16}O} रेखा गैस की मात्रा के समानुपाती बनी रहती है,
\ce{^{12}C^{16}O} रेखा ताप देती है: साथ मिलकर वे दोनों को मापती हैं।
\textbf{19.} $\tilde\nu_0 = 2169.81 - 2 \times 13.288 = \qty{2143.24}{cm^{-1}}$, \qty{4.666}{\micro m}।
\textbf{20.} $2\tilde\omega_e - 6\tilde\omega_ex_e = \qty{4259.89}{cm^{-1}}$।
\textbf{21.} $G(0) = \frac12\tilde\omega_e - \frac14\tilde\omega_ex_e = \qty{1081.58}{cm^{-1}}$।
\textbf{22.} $R(0) - P(1) = 2(B_0 + B_1) \approx 4B \approx \qty{7.7}{cm^{-1}}$।
\textbf{23.} \ce{^{13}C^{16}O}: भारी \omterm{def:b3:quantum-model-systems:oscillator}{समानीत द्रव्यमान}, निम्नतर आवृत्ति।
\textbf{24.} \textbf{$r_0(\ce{CO}) = \qty{113.1}{pm}$}, एक ही रेडियो आवृत्ति से।
\end{solution}
